% Generated draft: requires mathematical and source review.
\documentclass[11pt]{article}
\usepackage[margin=25mm]{geometry}
\usepackage{fontspec}
\setmainfont{Latin Modern Roman}
\usepackage{amsmath,amssymb,mathtools,unicode-math}
\setmathfont{Latin Modern Math}
\usepackage{enumitem,needspace,xurl,hyperref}
\setlength{\parindent}{0pt}
\setlength{\parskip}{5pt}
\setlength{\emergencystretch}{3em}
\newcommand{\N}{\mathbb{N}}
\newcommand{\Z}{\mathbb{Z}}
\newcommand{\Q}{\mathbb{Q}}
\newcommand{\R}{\mathbb{R}}
\newcommand{\C}{\mathbb{C}}
\newcommand{\F}{\mathbb{F}}
\newcommand{\A}{\mathbb{A}}
\newcommand{\PP}{\mathbb{P}}
\newcommand{\E}{\mathbb{E}}
\newcommand{\eps}{\varepsilon}
\newcommand{\dd}{\,\mathrm{d}}
\begin{document}
\section*{UnsolvedMath ID: 6800015}
\textbf{Source problem code:} catalogue entry
\subsection{Definitions and mathematical statement}
A Lie groupoid $\mathcal G\rightrightarrows M$ has smooth source, target, unit, inverse, and multiplication maps. A symplectic groupoid carries a closed nondegenerate two-form $\omega$ satisfying multiplicativity,
\[
m^*\omega=\operatorname{pr}_1^*\omega+\operatorname{pr}_2^*\omega
\quad\text{on the composable-pair space }\mathcal G^{(2)}.
\]
Its Lie algebroid is the source-vertical tangent bundle at the units, with anchor induced by the target map. In the usual Poisson correspondence, the algebroid of a symplectic groupoid is identified with a cotangent Poisson algebroid. This is not automatically the same assertion as putting a nondegenerate closed algebroid two-form on an arbitrary algebroid.

Explain the existence and role of the symplectic structure associated with the source's classifying groupoid/algebroid, and determine its relevance to moduli spaces of geometric structures of finite type. A complete formulation must identify the geometric-structure class, equivalence relation, groupoid and algebroid, and the proposed geometric structure on the moduli object.

The supplied item does not identify any of these referents or define finite type in this setting. It also does not specify whether the moduli space is a smooth quotient, a singular space, or a stack. Thus the task is a classification/Poisson-geometry research direction; the standard symplectic-groupoid definition above supplies context without asserting that every classifying groupoid admits such a form or that every Lie algebroid integrates.
\subsection{Short English statement}
Explain why the relevant classifying groupoid or algebroid should be symplectic and how that structure controls moduli of finite-type geometric structures. The actual objects, finite-type condition, and moduli category are unspecified.
\subsection{Sources}
\begin{itemize}
\item Morgan and Pansu, \emph{A List of Open Problems in Differential Geometry} (2018), Question environment 15, ``Classification problems and Poisson structures,'' proposed by Rui Loja Fernandes, University of Illinois, \url{https://www.imo.universite-paris-saclay.fr/~pansu/problems_MTDG.tex}; accessed 2026-07-29, source-TeX extraction.
\item The supplied research summary mentions Fernandes's Lie groupoid/algebroid work and Goldman-type moduli constructions without individual publication identifiers or precise theorems.
\item Local registry \texttt{data/unsolvedmath/problems.json}, permanent UnsolvedMath ID $6800015$, source code catalogue entry, \emph{Classification problems and Poisson structures}. Dataset attribution: Ulam AI and the UnsolvedMath Contributors.
\end{itemize}
\end{document}
